\section{Discussion}
\label{sec:discussion}

\paragraph{Evaluating latent channels.}
The donor control is cheap, one extra forward pass per donor message, and the choice of donor changes the conclusion. With donors from other documents, nearly all of the benefit of text, of the receiver's own states and of real-text bridges in the main pair is content; with donors from the same document, about half of that content is silently counted as generic. Work that contrasts matched with swapped messages, for training \citep{du2025interlat} or evaluation, should state where the swapped messages come from. We also recommend reporting the receiver's own states injected the same way. In our setting this reference shows how quickly the benefit of a message falls with injection depth, and that even at the shallowest depth it delivers less content than the corresponding text.

\paragraph{Building bridges without real corpus text.}
An LLM that wrote encyclopedia-style text without access to any corpus supplied probe text good enough for a bridge that delivers 82\% of the receiver-state reference's content. The models' own samples and weight algebra are much weaker substitutes. In the main WikiText experiment, the tested synthetic fitting recipes cost more than content: with these bridges a message about a different document leaves the receiver worse off on average than no message at all, while the real-text bridge and plain text do not. In other settings real-text bridges and plain text can also be hurt by irrelevant messages, and on reviews the self-sample bridge's generic part is not negative (\S\ref{sec:scale}); we compare complete recipes and do not isolate synthetic provenance as the cause. For systems that may receive irrelevant messages this is a robustness cost that a no-context comparison hides; behaviour under adversarially chosen messages remains untested. We do not know whether a small amount of in-domain text would remove it; testing that directly is an obvious next step. \SCALEDISCUSSION

\paragraph{Relation to weight-space reasoning.}
The study began from evidence that frontier models can read and edit network weights \citep{takehi2026sparks}. In small testbeds an LLM designed complete bridges from weights alone (Appendix~\ref{app:small}). At pretrained-model scale its useful contribution narrowed to supplying in-register text, and the distribution-matching recipe it had proposed on the small testbed; our experiments do not separate the LLM's world knowledge from register and topic coverage as the reason the text works. The weight-only map we tested, a least-squares fit between embedding matrices in the spirit of the within-model alignment of LatentMAS \citep{zou2025latentmas}, delivered 2\% of the reference. This rules out that particular map, not weight-based alignment in general.

\section{Limitations}
\label{sec:limitations}
\begin{itemize}
\item \textbf{Scale and family.} \XFAMILYLIMIT
\item \textbf{One task, few domains.} We measure next-token prediction, fitted on WikiText-like text and tested on WikiText, long movie reviews and Python code. We do not test whether the measured content predicts benefits in downstream collaboration such as question answering or reasoning.
\item \textbf{Cross-family protocol.} Cross-family pairs use a sender message length chosen from the test inputs, so the sender often reads less of the prefix than the receiver-state reference (93\% of its tokens for GPT-2 into Pythia, 74\% into Qwen2.5), and we have no position-matched no-context control for receivers with absolute positions. Their results are protocol-specific case studies. Two further planned pairs were dropped for compute after the first cross-family run had started (Appendix~\ref{app:general}), so the three reported pairs are not a blind or representative sample, and their intervals do not include pair-selection uncertainty.
\item \textbf{One receiver design.} The receiver gets a single residual-stream injection at one layer. Multi-layer or key--value injection may raise what latent messages can deliver, and the text-versus-latent comparison is specific to this design.
\item \textbf{Synthetic corpora.} Main results use one LLM, one prompt and 64 topic areas; the variant experiment adds one corpus from a second LLM and one fiction corpus, each a single generated sample that also differs in token count and topic coverage. The second LLM was run through an agentic command-line tool in a read-only sandbox, which restricts but does not remove tool access; we did not pin its model version. We cannot rule out that either LLM's pretraining included WikiText-like text.
\item \textbf{Log-score, not information.} Gains are differences in expected log-score through a fixed receiver. They measure usefulness to that receiver, not the mutual information between message and prefix.
\item \textbf{Statistics.} Intervals resample test articles, fitting seeds and donor draws, but with only two or three fits per condition the seed component is poorly estimated, and the donor pool itself is fixed (donors are other test windows). Per-seed values are reported for the main pair. Diagnostics (\S\ref{sec:diag}) and layer sweeps use one seed. \SCALESEEDS
\end{itemize}
